\section{讲座定位与问题意识}

本讲由 Sierra 联合创始人 Clay Bavor 主讲，主题是“把 Agent 真正部署到真实客户场景之后，会遇到哪些不在 Demo 里的硬问题”。这节课和偏研究的 Agent 课程不同，重点不在“能不能做出一个会推理的系统”，而在“能不能让这个系统在成本、质量、合规、稳定性四个维度同时站住”。

讲者开场提到，当 Sierra 在 2023 年初和客户讨论 Agent 时，很多人还不理解“Agent”是什么；30 个月后，这个方向已经进入 Berkeley 的研究生课程体系。这个变化速度本身就是课程背景：Agent 技术栈迭代非常快，产品和组织如果不能按季度重构，就会被现实流量与异常输入打穿。

\textit{``This will be a practical talk... making contact with reality again and again and again.''}

\begin{knowledgebox}{课程主线}
这节课围绕三条主线展开：第一，customer-facing agent 的业务价值如何定义；第二，语音、评测、观测、安全这些“冰山水下部分”如何工程化；第三，企业如何在 build、buy、build-with 三种策略间找到可持续路径。
\end{knowledgebox}

\begin{importantbox}{读这份笔记的方式}
建议按“业务目标 $\rightarrow$ 系统能力 $\rightarrow$ 运行保障”三层阅读：
\begin{itemize}[leftmargin=1.5em]
    \item 先看第 2--5 节，明确为什么客服 Agent 的核心指标是“可验证解决率”而不是“对话看起来很聪明”；
    \item 再看第 6--10 节，理解 memory、voice pipeline、simulation、tracing、guardrails；
    \item 最后看第 11--12 节，把技术方案映射到组织与路线图。
\end{itemize}
\end{importantbox}

\subsection{本章小结}

这节课本质是一门“Agent 生产工程学”。讲者反复强调：Demo 不是产品，产品不是系统图，系统图也不是可运营业务。后续章节将沿着这一逻辑，把“可演示能力”转换为“可经营能力”。

